\section{Experiments}
Useful graph adaptation requires three things: the simulator must be able to use additional messages, the controller must place them where they improve a prediction, and that improvement must persist as predictions feed back into the next state. We test these requirements in turn. The first comparison holds the expansion policy fixed and changes training. The second freezes each trained model and compares policies on identical observed histories. Full rollouts then expose the controller to the states it creates.

\subsection{Can the simulator learn to use extra messages?}
We pair base-only training with mixed-graph training, which independently expands each training example with probability $1/2$ by adding a uniformly sampled quarter of the optional annulus pairs. Both arms use the faithful objective and matched initialization, frames and noise. Goop and Sand train from initialization for 100k updates; WaterDrop uses paired 10k continuations of three inspected 100k parents, so its effects are conditional on those parents. Each comparison retains three paired seeds. The native capped base graph and self-messages are preserved; architecture and protocol details appear in Appendix~\ref{sec:implementation-details}.

Goop shows that training support matters even for an allocator with no learned ranking. Under the same random25 evaluation policy, full-rollout MSE falls from $0.301\pm0.169$ to $0.129\pm0.031$ after graph exposure, improving in every seed (Figure~\ref{fig:goop-exposure-placement}A). This establishes a benefit of exposure for a fixed expansion policy. It does not yet show that residual scores can select useful edges, or that expansion improves on the native base graph.


\subsection{Does residual risk place the budget well?}
To isolate placement, we evaluate each policy on the same observed test histories within each material: 150 per Goop or Sand model and 297 per WaterDrop endpoint. Goop and Sand compare random25, speed25, risk25 and RMS25 at the same quarter-annulus budget and radius ratio 1.267; base and dense control expansion itself. WaterDrop uses the same policy set except RMS25, which was not predeclared. Speed uses displacement magnitude, RMS uses relative velocity over native neighbors, and risk uses scores from the preceding observed base history. These comparisons hold the input state and model fixed.

The exposure gain does not resolve Goop's placement problem. Previous-observed risk remains worse than random in every mixed-model seed, although its mean deficit narrows from $(6.39\pm1.45)\times10^{-10}$ to $(2.84\pm3.62)\times10^{-10}$ (Figure~\ref{fig:goop-exposure-placement}B; Table~\ref{tab:goop-exposure-placement}). Here, realized action benefit is base-graph loss minus expanded-graph loss on the same history and target. The residual score correlates with base error at $.214\pm.088$, but with the benefit of its selected sparse action at only $.001\pm.048$. The score identifies difficult predictions much better than it identifies predictions that its chosen expansion improves.

WaterDrop tests whether exposure can change the placement ordering. It reverses the observed-test risk-minus-random ordering in every seed: the mean gap changes from $(+1.62\pm1.07)\times10^{-10}$ to $(-0.99\pm0.67)\times10^{-10}$. The validation gap also improves, but mixed risk remains worse than random there on average. We call the mixed-minus-base change in the risk-minus-random gap the \emph{training interaction}; a negative value means exposure narrows that gap. Sand shows a third pattern: exposure narrows its test deficit in every seed, from $(+4.76\pm3.80)\times10^{-9}$ to $(+1.87\pm1.04)\times10^{-9}$, while risk continues to lose to random in both arms. A favorable training interaction can therefore coexist with ineffective placement.

\begin{table}[t]
\centering\scriptsize
\caption{Training and placement effects differ across materials. Observed-test position MSE contrasts are means $\pm$ sample SD over three paired seeds. Column scales differ. Risk uses the preceding observed base graph; negative risk-minus-random favors risk. Goop and Sand use fresh 100k training, WaterDrop a paired continuation to 110k.}
\label{tab:goop-exposure-placement}
\setlength{\tabcolsep}{2.5pt}
\renewcommand{\arraystretch}{1.1}
\begin{tabular}{lccc}
\toprule
Contrast & \shortstack{Goop\\$\times10^{-10}$} & \shortstack{WaterDrop\\$\times10^{-10}$} & \shortstack{Sand\\$\times10^{-9}$}\\
\midrule
\multicolumn{4}{l}{\textit{Training effect: mixed $-$ base-only}}\\
Base evaluation & $-3.04\pm5.12$ & $+0.70\pm0.40$ & $-0.84\pm2.91$\\
Random25 evaluation & $-5.51\pm4.28$ & $+0.28\pm0.92$ & $-3.85\pm3.66$\\
\midrule
\multicolumn{4}{l}{\textit{Placement gap: risk $-$ random}}\\
Base-only training & $+6.39\pm1.45$ & $+1.62\pm1.07$ & $+4.76\pm3.80$\\
Mixed-graph training & $+2.84\pm3.62$ & $-0.99\pm0.67$ & $+1.87\pm1.04$\\
Change in gap & $-3.55\pm2.18$ & $-2.61\pm0.74$ & $-2.89\pm2.81$\\
\bottomrule
\end{tabular}
\end{table}


The separate Goop3D 25k observed comparison also narrows the test risk gap, while fixed-random exposure improves only one of three seeds (Appendix~\ref{sec:goop3d-25k-exposure}). Training support and placement remain separate questions at this endpoint.

\subsection{Do the placement findings survive feedback?}
An observed-history comparison removes a major source of variation, but a deployed simulator repeatedly predicts its own inputs. Its cached-risk controller also uses scores from its preceding selected graph, whereas the observed test scores a preceding base graph. We therefore evaluate all declared policies over each complete forecast horizon: H395 for Goop, H995 for WaterDrop and H314 for Sand. Frames average within trajectories, trajectories within seeds, and reported means and sample SDs retain all three seeds.

WaterDrop illustrates why the observed ordering alone is insufficient. All 810 required rollout outcomes complete, and exposure improves random, dense, speed and cached risk in every seed. Yet the autonomous risk-minus-random training interaction is $+.00284\pm.00466$, with mixed signs across seeds. The observed-test reversal does not yield a consistent autonomous allocation gain.

Sand is the stronger counterexample. All 1,080 outcomes complete, but exposure worsens cached-risk error in every seed ($+.00917\pm.01201$) and widens its risk-minus-random gap in every seed ($+.01975\pm.03062$). This happens despite the smaller observed-history deficit. Random improves in only two seeds, and native base beats random in every seed of both arms. Dense exposure improves every seed, but dense still loses to both native base and random. These comparisons separate a model's response to training from the value of expanding its graph.

Goop completes 1,077 of 1,080 required outcomes. Its full-horizon risk conclusion is limited by three candidate-pair guard failures in mixed seed 2, one each under base, dense and cached risk. The affected mixed-arm means and the cached-risk training interaction remain undefined. The two defined risk-minus-random seed comparisons have opposite signs. Among complete mixed-arm policies, speed25 and RMS25 beat random25 in every seed, with MSE $.112\pm.027$ and $.125\pm.032$, respectively (Table~\ref{tab:goop-rollouts}).

\begin{table}[t]
\centering\scriptsize
\caption{Graph exposure has different effects across rollout policies. Entries are mixed-minus-base full-horizon position MSE, mean $\pm$ sample SD over three paired seeds; negative favors mixed training. Goop H395 and Sand H314 use 90 outcomes per arm/policy; WaterDrop H995 uses 81. All outcomes complete except one candidate-pair guard each in mixed Goop base, dense and cached risk: $\dagger$ leaves those full-horizon effects undefined. NP: policy not predeclared.}
\label{tab:goop-rollouts}
\setlength{\tabcolsep}{3pt}
\renewcommand{\arraystretch}{1.2}
\begin{tabular}{lccc}
\toprule
Evaluation policy & Goop & WaterDrop & Sand\\
\midrule
Base & $\dagger$ & $-0.005\pm0.008$ & $-0.00004\pm0.003$\\
Dense & $\dagger$ & $-0.012\pm0.004$ & $-0.003\pm0.003$\\
Random25 & $-0.172\pm0.163$ & $-0.006\pm0.006$ & $-0.011\pm0.019$\\
Speed25 & $-0.160\pm0.137$ & $-0.005\pm0.006$ & $+0.010\pm0.019$\\
Cached risk25 & $\dagger$ & $-0.003\pm0.001$ & $+0.009\pm0.012$\\
RMS25 & $-0.208\pm0.204$ & NP & $-0.003\pm0.005$\\
\bottomrule
\end{tabular}
\end{table}


Lower rollout MSE also leaves substantial physical mismatch. Goop's mixed random, speed and RMS policies place 19.55--22.16\% of particles more than $10^{-6}$ outside the metadata box, compared with $0.0632\%$ for truth. These geometric diagnostics do not test conservation. Complete policy results, excursions and failure accounting are reported in Appendices~\ref{sec:goop-graph-exposure}, \ref{sec:waterdrop-110k-exposure} and~\ref{sec:sand-graph-exposure}; Figure~\ref{fig:qualitative-goop2d} shows a fixed metadata-selected example.

